\section{Nuclear engineering feasibility and safety}
\label{sec:nuclear-safety}

\subsection{Evidence maturity}
JAEA's 30 MWth HTTR is a constructed helium-cooled, graphite-moderated test reactor that achieved a 950$^\circ$C outlet temperature at full power~\cite{jaeahttr2026ops}. The selected 600 MWth GTHTR300C-class configuration is instead a design/economic architecture~\cite{nishihara2007potential}, while INL TEV-961 is a process model showing that 925$^\circ$C reactor outlet and 900$^\circ$C delivered heat can supply the 176.8 MWth reforming duty~\cite{inltev961}. The project therefore combines demonstrated temperature precedent, designed reactor architecture and modelled process integration; it does not claim a demonstrated 600 MWth nuclear-SMR plant.

\subsection{TRISO and passive/inherent characteristics}
\textbf{Plain-language takeaway:} each TRISO particle is a tiny fuel kernel wrapped in several engineered barrier layers; those layers strongly retain fission products, but they are not indestructible. TRISO fuel is a layered pressure vessel at particle scale. Around the fuel kernel, a porous carbon buffer accommodates fission gas and kernel swelling; the inner pyrolytic-carbon (IPyC, a dense carbon coating) layer provides structural support and protects the silicon-carbide deposition interface; dense silicon carbide (SiC, a ceramic barrier) is the principal pressure-bearing and metallic-fission-product barrier; and the outer pyrolytic-carbon (OPyC) layer provides additional mechanical/environmental protection~\cite{sawa2001httrfuel,sawaueta2004}. INL's AGR qualification programme adds modern UCO irradiation and post-irradiation safety evidence: 1600--1800$^\circ$C tests measured fission-product release and identified the consequences of SiC/coating failures~\cite{demkowicz2015facs,morris2016agr1,inlagrrelease}. Intact SiC showed excellent cesium retention in these tests, but measurable failure/release mechanisms remain; TRISO therefore reduces source-term potential (the amount of radioactive material available for release) rather than making radionuclide release impossible.

HTGR safety also uses reactor-scale physics. JAEA reports HTTR average power density of 2.5 MW/m$^3$ and a large graphite inventory~\cite{jaeahttr2026ops}. In the 9 MW HTTR loss-of-forced-cooling test, primary circulators were stopped without control-rod insertion; test/analysis evidence shows power falling rapidly toward decay heat through negative temperature-reactivity feedback while temperatures changed slowly because of graphite thermal inertia~\cite{takamatsu2011lofc}. These features reduce transient severity but do not remove the need for shutdown systems, decay-heat analysis, confinement, emergency planning or accident assessment.

\subsection{Accident mechanisms retained}
\textbf{Plain-language takeaway:} favourable HTGR physics does not remove accident analysis; loss of helium pressure, unwanted air or steam entry, loss of heat removal and radioactive-material transport still have to be evaluated. NRC HTGR/TRISO PIRTs explicitly retain depressurisation accidents with air ingress and water ingress, while the NGNP process-heat PIRT identifies intermediate-loop blowdown/loss of heat sink, leakage into the primary system and reactive-gas ingress as coupled-system safety phenomena~\cite{nrctriso6844,nrcngnppirt2010}. Depressurisation and subsequent air ingress can expose hot graphite to oxygen, so graphite oxidation and its effects on structural integrity remain accident mechanisms; fission-product transport during normal and accident conditions is itself a dedicated NRC PIRT topic~\cite{nrcngnpfpt2008}. Water/steam ingress, loss of heat sink and primary leakage likewise require design-specific analysis. The project therefore does not describe HTGRs as risk free or claim that severe fuel damage is impossible under every condition.

\subsection{IHX: heat-transfer equipment and safety boundary}
\textbf{Plain-language takeaway:} the IHX is both a heater and a physical separator; its metal boundary must survive high temperature for long periods without allowing the reactor and chemical loops to mix. JAEA's constructed HTTR IHX is a 10 MW helium--helium exchanger whose Hastelloy XR tubes form the pressure boundary between primary and secondary helium. Development work addressed creep collapse and creep-fatigue (slow high-temperature deformation combined with repeated thermal/mechanical cycling), seismic behaviour, bundle thermal hydraulics and in-service inspection~\cite{hamamoto2006ihx}. Structural analysis reports primary helium near 950$^\circ$C and secondary helium up to about 905$^\circ$C, with high-temperature creep explicitly governing design~\cite{httrihx1995}.

IAEA coupling guidance identifies the intermediate circuit as a means to separate the nuclear island from the chemical plant, limit radioactive contamination and reduce propagation of chemical disturbances into the reactor~\cite{iaeanuclearhydrogen2013}. This supports the project's primary-helium $\rightarrow$ IHX $\rightarrow$ secondary-helium $\rightarrow$ reformer architecture. It does not remove scale-up work: a constructed 10 MW IHX and a designed 370 MWth source heat branch are not a completed 176.8 MWth project exchanger.

\subsection{Coupled nuclear--chemical transients}
\textbf{Plain-language takeaway:} if one plant suddenly changes state, the other must not be driven into an unsafe condition; ``propagation'' here means one plant's disturbance causing a consequential event in the other. A chemical-plant trip or rapid heat-demand change alters secondary-helium return conditions. IAEA/JAEA work identifies return-helium control and prevention of reactor trips as explicit integration-development needs~\cite{iaeacoupling2009}. Conversely, reactor trip or loss of nuclear heat requires the chemical plant to reach a safe state with hot, pressurised methane/steam/hydrogen inventories. Independent trips, isolation, secondary-loop thermal management and process depressurisation therefore require future transient design; this screening study does not size those systems.

\subsection{Chemical/process hazards and co-location}
The hydrogen plant introduces combustible hydrogen and methane, toxic carbon monoxide, hot high-pressure syngas/steam, reformer fire/explosion hazards, solvent/capture equipment and compressed CO$_2$. A future process hazard analysis must evaluate leakage, ignition, explosion overpressure, toxic/asphyxiant releases, escalation and loss of utilities. The NRC NGNP process-heat/hydrogen PIRT emphasises chemical inventories, layout and separation distance as major controls against propagation to the nuclear plant~\cite{nrcngnppirt2010}; IAEA work likewise identifies combustible leakage/explosion as an HTGR-hydrogen integration issue~\cite{iaeacoupling2009}.

\subsection{Safety-feasibility conclusion}
At screening level the design has a credible physical safety architecture: TRISO fuel, graphite/helium HTGR characteristics, an intermediate helium boundary and demonstrated high-temperature IHX precedent. \textbf{Safety feasibility remains conditional}: site-specific accident consequences, chemical-to-nuclear propagation limits, separation distances, detailed IHX integrity, emergency planning and regulatory acceptance are not demonstrated here.